\documentclass[12pt,a4paper]{article}

\usepackage[a4paper,margin=1in]{geometry}
\usepackage[draft]{graphicx}
\usepackage{fancyhdr}
\usepackage{hyperref}
\usepackage{array}
\usepackage{longtable}
\usepackage{setspace}
\usepackage{titlesec}

\hypersetup{
    colorlinks=true,
    linkcolor=black,
    urlcolor=blue,
    citecolor=black,
    pdftitle={AI - Trip Planner Project Review},
    pdfauthor={K Druthendra, K Navyatha Simha, N Roopaswi}
}

\pagestyle{fancy}
\fancyhf{}
\lhead{ISWE302L - Artificial Intelligence}
\rhead{AI - Trip Planner}
\cfoot{\thepage}

\setstretch{1.2}
\setlength{\parindent}{0pt}
\setlength{\parskip}{6pt}

\begin{document}

\begin{titlepage}
    \centering
    \vspace*{1.5cm}
    {\Large \textbf{Vellore Institute of Technology (VIT), Vellore} \par}
    \vspace{0.5cm}
    {\large School of Computer Science Engineering and Information Systems \par}
    \vspace{1.2cm}

    {\Large \textbf{Project Review Report} \par}
    \vspace{0.8cm}

    {\large \textbf{Course Name:} ISWE302L -- Artificial Intelligence \par}
    \vspace{0.4cm}
    {\large \textbf{Semester:} Winter Semester 2025--26 \par}
    \vspace{0.8cm}

    {\LARGE \textbf{Project Title} \par}
    \vspace{0.3cm}
    {\Huge \textbf{AI -- Trip Planner} \par}
    \vspace{0.6cm}
    {\large \textbf{Domain:} Artificial Intelligence \& Mobile Application Development \par}

    \vspace{1.2cm}
    {\Large \textbf{Team Members} \par}
    \vspace{0.4cm}

    \renewcommand{\arraystretch}{1.3}
    \begin{tabular}{|>{\centering\arraybackslash}m{1.2cm}|m{8cm}|>{\centering\arraybackslash}m{3.5cm}|}
        \hline
        \textbf{S.No} & \textbf{Name} & \textbf{Register Number} \\
        \hline
        1 & K Druthendra & 23MIS0213 \\
        \hline
        2 & K Navyatha Simha & 23MIS0242 \\
        \hline
        3 & N Roopaswi & 23MIS0217 \\
        \hline
    \end{tabular}

    \vfill
    {\large Academic Year: 2025--2026 \par}
\end{titlepage}

\pagenumbering{roman}
\tableofcontents
\newpage
\listoffigures
\newpage

\pagenumbering{arabic}

\section{Abstract}
Travel planning is often time-consuming and complex, requiring users to manually search for destinations, hotels, transport, and daily schedules across multiple platforms. Most existing travel applications provide static recommendations based primarily on filters and user ratings. These systems generally lack personalization, intelligent itinerary generation, and real-time contextual decision-making. As a result, users frequently switch between different applications for booking, navigation, and trip management, which reduces both efficiency and overall user experience.

The proposed system, \textbf{AI Trip Planner}, is an intelligent mobile application built using React Native, Google Maps API, and Gemini AI to generate personalized travel itineraries. The application collects user preferences such as destination, number of days, budget, and interests. Using Gemini AI, it dynamically generates a complete travel plan that includes daily schedules, places to visit, optimal timings, hotel suggestions, and transport options. Integration with Google Maps enables real-time location visualization, route navigation, and distance estimation. The system also redirects users to official booking platforms for flights and hotels.

The major advantages of the proposed solution include automation of planning tasks, significant reduction in planning time, AI-driven personalization, and seamless integration of navigation and booking services.

\section{Keywords}
Artificial Intelligence, Trip Planning, Gemini AI, Google Maps API, React Native, Smart Tourism

\section{Introduction}
\subsection{Background}
The travel and tourism industry has undergone rapid digital transformation with the availability of online booking, navigation, and recommendation platforms. Despite these developments, end-to-end trip planning remains fragmented, requiring users to combine information from multiple sources.

\subsection{Problem Statement}
Current travel planning tools are often static and do not effectively adapt to individual user needs such as budget constraints, interests, and trip duration. The absence of intelligent itinerary generation and integrated services creates inefficiency in planning and execution.

\subsection{Objectives}
The key objectives of this project are:
\begin{itemize}
    \item To develop an intelligent mobile-based travel planner.
    \item To generate personalized itineraries using Gemini AI.
    \item To integrate route visualization and navigation via Google Maps API.
    \item To simplify travel organization through unified planning and booking redirection.
\end{itemize}

\section{Literature Review}
Recent studies in smart tourism emphasize intelligent decision support and context-aware personalization. Gretzel et al. highlighted the role of data-driven ecosystems in smart tourism. Li et al. discussed intelligent tourism information systems and adaptive travel planning. Zhang et al. demonstrated the significance of deep learning in recommendation systems for personalized outputs. Foundational AI work by Goodfellow et al. supports the use of neural models in prediction and recommendation workflows. Borras et al. reviewed recommender systems specifically for tourism applications. Industry documentation from Google Developers and practical frameworks such as React Native Community resources provide implementation foundations for geolocation-enabled mobile systems. Contemporary trends in AI-assisted planning and human-centric tourism technologies are further discussed by OpenAI \& DeepMind and by Xiang \& Fesenmaier.

\section{System Analysis}
\subsection{Existing System}
Existing travel applications largely provide search and recommendation modules with static filtering (cost, rating, popularity). Users still manually create itineraries and verify timing feasibility, resulting in fragmented planning.

\subsection{Proposed System}
The proposed AI Trip Planner accepts user preferences and dynamically generates complete day-wise itineraries. It combines AI-generated recommendations with map-based route intelligence and booking redirection for practical execution.

\subsection{Advantages}
\begin{itemize}
    \item Reduced planning time through automation.
    \item Personalized recommendations based on user profile and constraints.
    \item Integrated map visualization and navigation support.
    \item Improved usability by reducing dependence on multiple apps.
\end{itemize}

\section{System Design}
\subsection{Architecture}
Figure~\ref{fig:architecture} presents the high-level architecture of the AI Trip Planner system.

\begin{figure}[h!]
    \centering
    \includegraphics[width=\textwidth]{architecture.png}
    \caption{System Architecture of AI -- Trip Planner}
    \label{fig:architecture}
\end{figure}

\subsection{Workflow Explanation}
The workflow begins with user input collection through the mobile interface. The input payload is processed by the trip generation module, which invokes Gemini AI to produce personalized itinerary content. Simultaneously, Google Maps services provide geographic context, route options, and distance estimates. The generated plan is displayed in structured day-wise format, along with hotel and transport suggestions. Users can then navigate to official external booking portals for reservation completion.

\section{Implementation}
\subsection{React Native}
React Native is used to build a cross-platform mobile interface for Android-based deployment and testing. It supports modular UI development and efficient integration with third-party services.

\subsection{Gemini API}
Gemini API is used as the intelligence layer for itinerary generation. User constraints (destination, days, budget, interests) are transformed into prompt structures to produce meaningful travel plans.

\subsection{Google Maps API}
Google Maps API is integrated to display locations, routes, and estimated distances. It enhances user understanding of travel feasibility and place connectivity.

\subsection{.env Configuration}
Sensitive credentials are managed using environment configuration.

\begin{verbatim}
GEMINI_API_KEY=your_gemini_api_key
GOOGLE_MAPS_API_KEY=your_google_maps_api_key
\end{verbatim}

\section{Dataset Description}
The project utilizes the following datasets and data streams:
\begin{itemize}
    \item User input dataset (destination, budget, days, interests)
    \item Gemini AI generated dataset (itinerary, hotels, schedule)
    \item Google Maps dataset (coordinates, routes)
    \item Booking redirection links
\end{itemize}

\section{Hardware Requirements}
\begin{table}[h!]
    \centering
    \renewcommand{\arraystretch}{1.25}
    \begin{tabular}{|m{0.8cm}|m{10cm}|m{4cm}|}
        \hline
        \textbf{S.No} & \textbf{Requirement} & \textbf{Specification} \\
        \hline
        1 & Processor & Intel i5 or above \\
        \hline
        2 & Memory (RAM) & 8GB RAM \\
        \hline
        3 & Storage & 20GB storage \\
        \hline
        4 & Device & Android device/emulator \\
        \hline
        5 & Connectivity & Internet \\
        \hline
    \end{tabular}
    \caption{Hardware Requirements}
\end{table}

\section{Software Requirements}
\begin{table}[h!]
    \centering
    \renewcommand{\arraystretch}{1.25}
    \begin{tabular}{|m{0.8cm}|m{14.2cm}|}
        \hline
        \textbf{S.No} & \textbf{Software Requirement} \\
        \hline
        1 & React Native \\
        \hline
        2 & Node.js \\
        \hline
        3 & Android Studio \\
        \hline
        4 & VS Code \\
        \hline
        5 & Gemini API Key \\
        \hline
        6 & Google Maps API Key \\
        \hline
        7 & Firebase (optional) \\
        \hline
    \end{tabular}
    \caption{Software Requirements}
\end{table}

\section{Features}
\begin{itemize}
    \item AI itinerary
    \item Map integration
    \item Customization
    \item Real-time suggestions
\end{itemize}

\section{Screenshots Section}
\begin{figure}[h!]
    \centering
    \includegraphics[width=\textwidth]{screenshot1.png}
    \caption{Home Screen}
\end{figure}

\begin{figure}[h!]
    \centering
    \includegraphics[width=\textwidth]{screenshot2.png}
    \caption{User Preference Input}
\end{figure}

\begin{figure}[h!]
    \centering
    \includegraphics[width=\textwidth]{screenshot3.png}
    \caption{Destination Selection}
\end{figure}

\begin{figure}[h!]
    \centering
    \includegraphics[width=\textwidth]{screenshot4.png}
    \caption{Generated Itinerary View}
\end{figure}

\begin{figure}[h!]
    \centering
    \includegraphics[width=\textwidth]{screenshot5.png}
    \caption{Daily Plan Breakdown}
\end{figure}

\begin{figure}[h!]
    \centering
    \includegraphics[width=\textwidth]{screenshot6.png}
    \caption{Map and Route Visualization}
\end{figure}

\begin{figure}[h!]
    \centering
    \includegraphics[width=\textwidth]{screenshot7.png}
    \caption{Hotel Suggestions}
\end{figure}

\begin{figure}[h!]
    \centering
    \includegraphics[width=\textwidth]{screenshot8.png}
    \caption{Transport Recommendation Screen}
\end{figure}

\begin{figure}[h!]
    \centering
    \includegraphics[width=\textwidth]{screenshot9.png}
    \caption{Booking Redirection Interface}
\end{figure}

\begin{figure}[h!]
    \centering
    \includegraphics[width=\textwidth]{screenshot10.png}
    \caption{Final Trip Summary}
\end{figure}

\clearpage
\section{Results and Discussion}
The implemented AI Trip Planner successfully demonstrates intelligent itinerary generation with strong personalization capability. The integration of Gemini AI enables generation of practical and structured plans based on user constraints. Google Maps integration improves route awareness and usability. Overall, the system reduces manual effort and improves trip planning efficiency.

\section{Applications}
\begin{itemize}
    \item Personalized travel planning for individual users.
    \item Smart tourism assistance for domestic and international trips.
    \item Educational demonstration of AI integration in mobile apps.
    \item Prototype foundation for commercial travel recommendation platforms.
\end{itemize}

\section{Limitations}
\begin{itemize}
    \item Dependence on internet connectivity for AI and maps.
    \item AI responses may require occasional user validation.
    \item Limited offline functionality in the current version.
    \item Booking flow depends on external platforms and APIs.
\end{itemize}

\section{Future Enhancements}
\begin{itemize}
    \item Budget prediction
    \item AI chatbot
    \item Offline maps
    \item Group planning
\end{itemize}

\section{Conclusion}
AI -- Trip Planner presents an effective approach to modern travel organization by combining artificial intelligence, mobile application design, and geospatial services. The system addresses critical gaps in existing travel tools by automating itinerary generation and improving personalization. With further enhancements, the platform can evolve into a comprehensive intelligent travel assistant suitable for real-world deployment in smart tourism ecosystems.

\section{References}
\begin{enumerate}
    \setcounter{enumi}{19}
    \item Gretzel et al. (2015) Smart tourism
    \item Li et al. (2017) Tourism systems
    \item Zhang et al. (2020) Deep learning recommendation
    \item Goodfellow et al. (2016) Deep learning
    \item Borras et al. (2014) Recommender systems
    \item Google Developers (2023)
    \item OpenAI \& DeepMind (2024)
    \item Xiang \& Fesenmaier (2017)
    \item React Native Community (2024)
    \item Buhalis \& Amaranggana (2015)
\end{enumerate}

\end{document}
